%!TEX root = ../main.tex
\chapter{Introduzione}
\label{cap:introduzione}

Un incrocio semaforizzato stabilisce periodicamente quali flussi veicolari 
autorizzare e quali arrestare. Su una singola intersezione l'ottimizzazione 
del numero di veicoli che riescono ad attraversarla è banale:
si concede la precedenza alla direzione con il maggior traffico in attesa.
Tuttavia se si considera una rete composta da vari incroci e se si cerca inoltre di ottimizzare
anche il tempo di attesa dei veicoli, il problema diventa molto più complesso.
La fase selezionata da un incrocio influenza con un certo ritardo temporale il volume
di veicoli in arrivo alle intersezioni adiacenti. 
Pertanto il controllo semaforico su rete si configura 
come un problema di coordinamento distribuito ancor prima che di ottimizzazione isolata.

Attualmente la stragrande maggioranza delle intersezioni è regolata da sistemi a tempo fisso.
Tale approccio prevede una sequenza di fasi di durata predeterminata, calcolata una tantum sulla base di stime di traffico medie.
Il sistema non reagisce quindi in alcun modo alle condizioni osservate in tempo reale.
Al giorno d'oggi si possono comunque individuare alcuni casi in cui la scelta della fase semaforica viene influenzata dal traffico, 
come negli incroci che presentano un sensore di rilevamento di veicoli in prossimità del semaforo, grazie al quale
per esempio si attiva il flusso sulle strade secondarie solo quando si presenta fisicamente un veicolo.
Questi sistemi sono però molto limitati, sia dal punto di vista dei dati raccolti, sia dal punto di vista delle azioni eseguibili. 
La ricerca si è quindi concentrata verso lo studio di sistemi di controllo semaforico che potessero, almeno in parte, 
reagire più concretamente alle condizioni di traffico in tempo reale, necessitando però di una raccolta di dati più elevata.
Al momento si possono trovare nel mondo e in Italia incroci che sfruttano i risultati di queste ricerche, 
ma la loro diffusione non è ancora ampia a causa di ostacoli economici, tecnici, burocratici e normativi.

Detto ciò, lo studio di questi approcci reattivi rimane ad ogni modo interessante e di potenziale beneficio.
Uno degli algoritmi che hanno ottenuto maggiore popolarità in questo ambito è
\textsc{Max-Pressure} \citep{varaiya2013maxpressure}, il quale introduce una regola locale che seleziona a ogni istante la
fase con la pressione istantanea più alta (differenza tra veicoli in
ingresso e in uscita sulle corsie servite da quella fase). \textsc{Max-Pressure}
garantisce proprietà di stabilità della rete nel suo insieme, ma la sua
natura puramente greedy non offre alcuna garanzia di equità (\emph{fairness}):
in letteratura sul controllo semaforico, l'equità indica che tutti gli
utenti della strada ricevano un accesso equo alle risorse dell'intersezione,
evitando che alcuni veicoli o alcune direzioni sperimentino attese
eccessive rispetto ad altri \citep{xiao2025tscreview} — un obiettivo
concettualmente distinto dall'efficienza aggregata di rete (tempo di
percorrenza medio, throughput): un controllore può massimizzare il flusso
complessivo e allo stesso tempo lasciare una minoranza di veicoli
sistematicamente penalizzata, se la sua regola di decisione non include
alcun meccanismo che penalizzi esplicitamente le attese prolungate. È
esattamente il caso di \textsc{Max-Pressure}: un'intersezione posta su
un'arteria ad alto scorrimento tende a favorire sistematicamente la
direzione dominante, costringendo i flussi minoritari ad attese anche
indefinitamente lunghe, poiché la regola sceglie sempre la fase a pressione
istantanea più alta senza mai considerare da quanto tempo attende la
direzione esclusa — una criticità che questo lavoro osserva empiricamente
(Capitolo~\ref{cap:tsc-modelli}).

Alcuni metodi recenti in letteratura affrontano il problema penalizzando
esplicitamente le fasi responsabili di attese eccessive, tipicamente solo
oltre una soglia, per non sacrificare le prestazioni medie quando non è
necessario: FELight, ad esempio, introduce una metrica di equità basata sui
tempi di attesa dei veicoli e penalizza selettivamente le fasi con ritardi
eccessivi, riportando una riduzione delle attese prolungate su direzioni
specifiche senza compromettere l'efficienza complessiva della rete
\citep{xiao2025tscreview}. Questa tesi persegue lo stesso obiettivo con un
meccanismo diverso, formalizzato nel Capitolo~\ref{cap:tsc-modelli}: un
termine dedicato della ricompensa penalizza l'attesa massima osservata su
una corsia non servita dalla fase corrente, pesato da un iperparametro che
ne modula l'intensità — concettualmente lo stesso principio (penalizzare la
coda peggiore, non solo ottimizzare il throughput aggregato), applicato però
a ogni passo di decisione per ogni intersezione, non a un sottoinsieme di
fasi identificate a priori. È il motivo per cui questo lavoro accetta, entro
certi limiti, un tempo di percorrenza medio peggiore pur di ottenere
un'attesa massima più contenuta: a parità di rete, un controllore che fa
aspettare pochi veicoli molto a lungo non è preferibile a uno leggermente
più lento in media ma che non lascia mai nessuno bloccato a tempo
indefinito.

Per misurare quanto questo meccanismo funzioni nella pratica, questo lavoro
non si affida al solo segnale di ricompensa — che dipende dai propri stessi
parametri e non è quindi un metro di giudizio imparziale — ma a una metrica
puramente osservazionale, calcolata identicamente per ogni controllore
indipendentemente da come è stato addestrato o se non è stato addestrato
affatto: la massima attesa storica \emph{consecutiva} osservata su una
corsia, in una qualunque intersezione della rete, $\text{wait}_{\max}$ —
cioè la durata del più lungo singolo stallo mai osservato, non un'attesa
cumulativa sull'intero viaggio del veicolo. La metrica include
deliberatamente gli stalli ancora in corso al termine dell'episodio, per lo
stesso principio di non-esclusione del caso peggiore che vale anche per il
tempo di percorrenza (definito formalmente nel Capitolo~\ref{cap:confronto},
§\ref{sec:metriche}): un veicolo fermo nello stesso punto dall'inizio del
blocco fino all'ultimo istante simulato non farebbe mai scattare un nuovo
record dopo la fine dell'episodio, e ignorare questi casi premierebbe
indirettamente chi lascia un veicolo bloccato per sempre invece di
penalizzarlo come merita. Il caso non è ipotetico: un'attesa di 1125
secondi registrata da \textsc{Max-Pressure} in uno degli esperimenti di
questo lavoro è stata verificata ricostruendo la traiettoria del veicolo
responsabile direttamente dal replay grezzo della simulazione — il veicolo
risultava fermo nello stesso identico punto per oltre mille secondi, quasi
l'intera durata dell'episodio simulato.

L'obiettivo di questo lavoro è dunque duplice, e i suoi due termini non
sono sempre allineati: da un lato ridurre il tempo di percorrenza medio e
massimo dei veicoli (efficienza), dall'altro impedire che una direzione di
traffico venga sistematicamente sacrificata a vantaggio delle altre
(equità). Il Capitolo~\ref{cap:confronto} mostra come questo compromesso si
comporti concretamente al variare del peso assegnato alla componente di
equità nella ricompensa.

A questi due obiettivi se ne affianca un terzo, di natura diversa dai
primi due: verificare che le prestazioni misurate non dipendano dalla
specifica rete stradale su cui il modello è stato addestrato, ma si
mantengano, almeno in parte, su topologie e condizioni di traffico mai
incontrate durante il training.

Gran parte della letteratura recente sul controllo semaforico appreso non
affronta sistematicamente questa domanda: la prassi più diffusa è
addestrare e valutare un modello sulla stessa rete stradale, al più
ripetendo l'intero esperimento, addestramento incluso, su una rete
diversa — non trasferendo un singolo modello già addestrato a una rete che
non ha mai visto. CoLight \citep{wei2019colight} conduce cinque esperimenti
indipendenti, due su reti sintetiche e tre su reti reali (Manhattan,
Hangzhou, Jinan), ciascuno con il proprio addestramento; MetaSTGAT
\citep{wang2022metastgat}, il modello di riferimento di questo lavoro,
segue lo stesso schema su una griglia sintetica e due dataset reali. Quando
in questi lavori compare il termine \emph{generalization}, indica quasi
sempre la tenuta a diverse intensità di traffico sulla stessa rete, non a
una topologia strutturalmente diversa \citep{wang2026astgcndrl,
  sun2026tgmaddpg}. Zhang et al. lo rendono esplicito nella motivazione di
GeneraLight \citep{zhang2020generalight}: la maggior parte dei modelli RL
per il controllo semaforico è addestrata e valutata nello stesso ambiente
di traffico, il che porta a overfitting e a una scarsa capacità di
generalizzazione quando le condizioni reali cambiano — un problema
riconosciuto abbastanza da giustificare un lavoro dedicato a risolverlo,
il che è di per sé indicativo di quanto la generalizzazione resti
l'eccezione, non la norma.

Le eccezioni che esistono restano circoscritte a un'unica dimensione di
variazione. AttendLight \citep{oroojlooy2020attendlight} addestra su un
insieme eterogeneo di intersezioni singole e ne verifica le prestazioni su
intersezioni mai viste, ma dichiara che l'estensione a una rete di più
intersezioni coordinate resta lavoro futuro. Advanced-XLight
\citep{zhang2022advancedxlight} misura un rapporto di trasferimento tra un
modello addestrato su un dataset reale e valutato direttamente su un
altro, senza ulteriore addestramento, ma sempre tra reti di scala
paragonabile. Lo stesso GeneraLight generalizza al variare del flusso di
traffico, non della topologia della rete.

Questo lavoro adotta invece la generalizzazione come requisito sistematico
su un asse che la letteratura esaminata non copre: ogni modello, RL o
baseline, è addestrato una sola volta su una griglia $4\times4$ e valutato,
senza alcun riaddestramento, anche su topologie di dimensione maggiore mai
incontrate (griglie $5\times5$ e $6\times6$) e su una spaziatura stradale
diversa, oltre che su profili di traffico differenti da quello di
addestramento (Capitolo~\ref{cap:confronto}, §\ref{sec:protocollo}).

Il controllo appreso tramite \emph{reinforcement learning} (RL) propone
un paradigma differente a quello di \textsc{Max-Pressure}: invece di una regola scritta a mano, una politica
$\pi$ viene addestrata a massimizzare una ricompensa che codifica l'obiettivo
di rete (tipicamente throughput e tempo di viaggio) osservando ripetutamente
le conseguenze delle proprie azioni in simulazione. Applicato a una singola intersezione, 
il problema è modellabile come un MDP standard; su scala di rete, assume la forma 
di un sistema multi-agente a osservabilità parziale, dove ogni nodo rileva 
esclusivamente il proprio intorno locale. Emerge quindi la necessità di un 
meccanismo per la condivisione delle informazioni tra agenti adiacenti. 
Le \emph{graph neural network} (GNN) assolvono a questa funzione: rappresentando la rete stradale
come un grafo, con le intersezioni come nodi e le strade come archi, 
l'informazione si propaga tra nodi vicini, permettendo a ogni intersezione di condizionare la propria
decisione sullo stato dei suoi vicini immediati (e non) senza richiedere una politica
centralizzata.

\textbf{MetaSTGAT} \citep{wang2022metastgat} combina questi aspetti
con uno ulteriore: invece di condividere gli stessi pesi tra tutte le
intersezioni della rete, un modulo di \emph{meta-learning} genera dinamicamente
i pesi di uno strato di attenzione e di uno strato ricorrente a partire
da caratteristiche locali dell'intersezione, così che nodi diversi possano essere serviti da
trasformazioni diverse pur restando un unico modello addestrato end-to-end.
Gli autori originali riportano una riduzione del travel time medio
rispetto a CoLight \citep{wei2019colight}, un metodo di riferimento 
basato su \emph{graph attention} ma con pesi condivisi. 
La combinazione di \emph{attention} su grafo, ricorrenza e reinforcement learning non è isolata: 
varianti di questo framework dominano la letteratura recente
\citep{wang2026astgcndrl,sun2026tgmaddpg,hu2026mpnnlight}.

\section{Contributo di questo lavoro}

Questa tesi replica da zero l'architettura e la procedura di addestramento di
MetaSTGAT, conducendo uno studio ablazionale sistematico e confrontandone le 
prestazioni con le baseline classiche del controllo semaforico. 
Nello specifico, i contributi sono:

\begin{itemize}
  \item un ambiente di simulazione multi-intersezione costruito su
    \textsc{CityFlow} \citep{zhang2019cityflow}, con reti stradali a griglia
    di dimensione $4\times4$, $5\times5$ e $6\times6$, otto fasi semaforiche
    per intersezione e una formulazione MDP dettagliata nel
    Capitolo~\ref{cap:tsc-modelli};
  \item una generazione parametrica dei dati di traffico, con due livelli di
    densità (\emph{flat} e \emph{peak}); in ciascuna
    configurazione viene potenziato artificialmente
    il traffico su una riga della griglia simulando un'arteria stradale, 
    in modo da verificare sistematicamente se un modello di controllo produce uno squilibrio
    direzionale quando il traffico non è simmetrico;
  \item un sistema di ablation che isola singolarmente ogni
    differenza tra l'implementazione proposta in questa tesi e la procedura descritta nell'articolo
    originale, organizzato in sei configurazioni che vanno dalla replica
    fedele del paper al modello più avanzato di questo lavoro;
  \item un protocollo di valutazione che riporta media e caso peggiore del tempo di viaggio
    dei veicoli, oltre alla massima attesa registrata da parte di un veicolo all'interno di un incrocio
\end{itemize}

Al momento della stesura di questo capitolo, l'estensione dell'architettura
verso meccanismi spaziali alternativi alla \emph{graph attention} (una
convoluzione grafica standard, GCN, e un meccanismo di propagazione a lungo
raggio ispirato a SONAR) è ancora in fase di progettazione e non è inclusa
nei risultati qui presentati. Il Capitolo~\ref{cap:tsc-modelli} ne introduce
comunque la motivazione, e il confronto sperimentale a parità di resto
dell'architettura è rimandato a un aggiornamento successivo di questo lavoro.

\section{Struttura della tesi}

Il Capitolo~\ref{cap:background} introduce il background necessario a
leggere i capitoli successivi senza assumere familiarità pregressa con
nessuno dei tre campi coinvolti: reinforcement learning, graph neural
network e controllo semaforico. Il
Capitolo~\ref{cap:tsc-modelli} formalizza il problema affrontato come MDP,
descrive l'ambiente di simulazione e le baseline classiche, e presenta
l'architettura di MetaSTGAT insieme a ogni differenza introdotta rispetto all'articolo 
originale, motivata caso per caso. Il Capitolo~\ref{cap:confronto}
descrive il protocollo sperimentale e riporta i risultati ottenuti,
incluso lo studio di ablation. Il Capitolo~\ref{cap:conclusioni} discute i
risultati alla luce della domanda di ricerca, dichiara i limiti dello studio
e indica il lavoro rimasto per la seconda parte del progetto.
